%!TEX program = lualatex
\documentclass[ngerman,headheight=70pt]{scrartcl}
\usepackage[ngerman]{babel}
\usepackage{blindtext}
\usepackage[utf8]{luainputenc}
\usepackage{scrlayer-scrpage}
\usepackage[T1]{fontenc}
\usepackage{fontspec}
\usepackage{multicol}
\usepackage{enumitem}
\usepackage[german=quotes]{csquotes}
\usepackage[bottom=2cm,footskip=8mm]{geometry}
\setmainfont{Source Sans Pro}
\MakeOuterQuote{"}

\setlength{\parindent}{0pt}
\setlength{\parskip}{6pt}

\pagestyle{plain.scrheadings}
\thispagestyle{scrheadings}
\lohead{\vspace*{-1.5cm}\includegraphics[scale=0.4]{../images/stupa_logo_rot.pdf}}
\rohead{\\Präsidium des\\Studierendenparlaments}

\lofoot[Universität Hamburg · Präsidium des Studierendenparlaments\\\hrule Von-Melle-Park 5 · D-20146 Hamburg]{Universität Hamburg · Präsidium des Studierendenparlaments\\\hrule Von-Melle-Park 5 · D-20146 Hamburg}
\cofoot[]{}
\cohead[\thepage]{}

\renewcommand{\thesubsection}{TOP \arabic{subsection}}
\renewcommand{\thesubsubsection}{\arabic{subsubsection}.}
\setcounter{subsection}{-1}

\begin{document}
    UHH · StuPa-Präsidium · Von-Melle-Park 5 · D-20146 Hamburg

    %%% PAGE ONE (EDIT LATER) %%%%
    % Hier den Titel des Protokolls eintragen
    \section*{Protokoll der 11. Sitzung des Studierendenparlaments vom 1. Dezember 2016}

    % Protokollanten eintragen
    \textbf{Protokoll: Gunhild Berdal}\\
    \textbf{Ort: VMP9 S08}\\
    \textbf{Beginn: 18. Uhr}\\
    \textbf{Ende: 23 Uhr}

    % Wohl die komplizierteste Stelle im Protokoll
    % Als Faustregel sollten nicht mehr als 3 Namen auf einer Zeile stehen
    % Einfach an dieser Vorlage orientieren
    % Im Idealfall ist hier nicht viel zu ändern, da kaum jemand fehlt.
    % Rücktritte müssen natürlich pro Protokoll angepasst werden
    \vspace{0.5cm}
    \begin{tabular}{ll}
        Anwesend: & \\
            RCDS (5 Sitze): & Kolja Kolb, Ramin Shakiba \\
             CampusGrün (14 Sitze): & Laura Franzen, Geoffrey Youett, Philipp Droll, \\
                                   & Yasemin Günther, Melf Johannsen, Mario Moldenhauer, \\
                                   & Christian Lagod, Jim Martens, Svenja Horn, \\
                                   & Martin Sievert, Armin Günther \\
             BGZM-Liste (2 Sitze) : & \\
             WiWi (2 Sitze): & \\
             Unicorns (5 Sitze): &  Johannes Peplow, Andreas Hartkamp \\
             Liste LINKS (3 Sitze): & Gunhild Berdal, Till Petersen, Sinah Mielich \\
             HWP (2 Sitze): &  \\
             MIN (4 Sitze): &  \\
             SDS* (3 Sitze): & Mena Winkler, Jacob Petersein, Artur Brückmann \\
             Bart-LISTE (2 Sitze): & Timo Zeimet \\
             LHG (1 Sitz): &  \\
             harte zeiten (1 Sitz): & Tobias Berking \\
             Jura (1 Sitz): & \\
             AL (2 Sitze): & Karima Schulze, Henri Weber \\
            & \\
        Entschuldigt: & Ramon Weilinger, Benjamin Welling (RCDS), Elena Rysikova, \\
                      & Ines Friedrich, Dean Haltmeyer (CG), Jakob Pape (BGZM), \\
                      & Charleen Lorenz (WiWi), Annkathrin Löffler, Marielle Hermstrüwer (Unicorns), \\
                      & Ailina Salten (MIN), Daniel Runge (BART)\\
                      &\\
        Unentschuldigt abwesend: & Jennifer Maack (RCDS), Karen Martirosian (BGZM), Claas-Friso Hente (WiWi), \\
                                & Katharina Kucza, Hakan Tasdemir, Gökce Celik (HWP),\\
                                & Lotte Rullkötter, Nabil Jabareen, Jan Detampel (MIN), \\
                                & Tobias Heisig (LHG), Maurice Martens (Jura) \\
                                &\\
        Rücktritte: & Antonia Niecke \rightarrow Kolja Kolb (RCDS) \\
                    & Tahnee Herzig \rightarrow Malte \rightarrow Ines Friedrich (CG) \\
                    & Dominic Laumer \rightarrow Daniel Runge (BART) \\
        Mandatsverluste: &  \\
    \end{tabular}
    \newpage
    %%% Vorgeschlagene Tagesordnung %%%
    \underline{Vorgeschlagene Tagesordnung}
    \begin{enumerate}[label={\textbf{Top \theenumi}},leftmargin=*]
        \item Aktueller Stand StuPa-Wahl (V1617-045) (30 Min.)
        \item Tag gegen Rechts (V1617-043, V1617-044) (45 Min.)
        \item Verschiedenes (5 Min.)
    \end{enumerate}

    \newpage

    % Hier beginnt Teil A. Bei Fortsetzungen kommt hier nur der Teil über
    % die Feststellung der endgültigen Fassung der Tagesordnung.
    % Für sonstige Sitzungen ist dies bereits das Gerüst für die Formalia.
    % Man muss nur noch die protokollierten Inhalte eintragen.

    {\Large \underline{Teil A}}

    \subsection{Formalia}



    \subsubsection{Geschäftsbericht Präsidium}


    \subsubsection{Anfragen an das Präsidium}


    \subsubsection{Geschäftsbericht AStA}



    \subsubsection{Anfragen an den AStA}



    \subsubsection{Dringlichkeitsanträge des AStA}

    Keine.

    \subsubsection{Aktuelle Stunde zum US-Wahlergebnis}

    \subsubsection{Feststellung der endgültigen Fassung des Teils B der Tagesordnung}


    \newpage
    %%% BESCHLOSSENE TAGESORDNUNG %%%
    \underline{Endgültige Fassung der Tagesordnung}
    \begin{enumerate}[label={\textbf{Top \theenumi}},leftmargin=*]
        \item Aktueller Stand StuPa-Wahl (V1617-045) (30 Min.)
        \item Tag gegen Rechts (V1617-043, V1617-044) (45 Min.)
        \item Verschiedenes (5 Min.)
    \end{enumerate}

    \subsubsection{Feststellung der Beschlussfähigkeit}


    \subsubsection{Genehmigung der Protokolle der vorangegangenen Sitzungen}

    \vspace{0.5cm}
    {\Large \underline{Teil B}}



\end{document}
